\clearpage
\SourceLeaf{83}
\Needspace{10\baselineskip}\begin{center}
\textbf{NEW PROOF OF LAGRANGE'S THEOREM.}
\end{center}

\noindent\hspace*{\parindent}The important theorem of Mr. \textsc{de la Grange}, according to which every function of
$x$ is represented by a series in powers of $u$ when $x=t+uX$,
was already proved ingeniously by its discoverer\HistoricalNote{*)}{*)}{Hist.\ de l'Ac.\ de Berlin T.\ XXIV, Année 1768 (Berlin 1770), p.\ 275. The entire paper can be found in translation in the third volume with which Mr.\ \textsc{Michelsen} supplemented \textsc{Euler}'s Analysis des Unendlichen. Recently Mr.\ D.\ \textsc{Murhard} also sought to make this proof better known by providing a concise extract in a program of his own.}.
This proof assumes the following form for $X$.
\begin{center}
$\displaystyle ax^2+bx^3+cx^4+.\,.\,.\,.$\HistoricalNote{**)}{**)}{The same applies to the inversion formula that Mr.\ \textsc{Eschenbach} discovered independently and Mr.\ M.\ \textsc{Rothe} proved, and which can easily be connected with \textsc{Lagrange}'s theorem.}
\end{center}
\noindent The proof is also rather lengthy. For these reasons, several geometers have sought
to derive the theorem solely from the nature of functions\HistoricalNote{***)}{***)}{\textsc{Lexell}, Comm.\ nov.\ Petr., T.\ XVI (1772), p.\ 230.\par
\textsc{Condorcet}, Miscell.\ Taurin., T.\ V (1770--73), Classe Math.\ p.\ 7.\par
\textsc{Cousin}, Astronomie physique, p.\ 15 (Paris 1787).\par
\textsc{Pfaff}, Archiv der r.\ u.\ a.\ M., 1.\ Band (1795), S.\ 81.}.
Of all these attempts, however, only the proof by Prof.\ \textsc{Pfaff} is completely
rigorous. The others are essentially mere induction. I hope some readers will
be pleased to find another equally rigorous proof here, which has nothing in
common with the proof by Prof.\ \textsc{Pfaff}. The main idea is essentially
the same as in \textsc{Cousin}'s proof, but the execution seems to me new and
simple. A variety of methods always serves to shed more light on a truth, and
the present theorem seems so far to be less well known among us than it deserves.

